\begin{table}[htbp]
\centering
\caption{Item analysis (classical test theory).}
\label{tab:items}
\small
\begin{tabular}{p{8.5cm}r}
\toprule
Measure & Value \\
\midrule
Items analysed & 166 \\
Difficulty (first-attempt \% correct), median [IQR] & 85.7 [75.3, 94.5] \\
Items with $<$50\% / 50--90\% / $>$90\% correct & 4 / 102 / 60 \\
Discrimination (item--rest $r$), median [IQR], items with $n\geq$30 & 0.22 [0.14, 0.32] ($k$ = 161) \\
Items with $r<0.20$ / $0.20$--$0.29$ / $\geq0.30$ & 71 / 42 / 48 \\
Share of errors on the most frequent distractor, median [IQR] & 71.8 [58.4, 84.3] \\
Observed distractors chosen by $<$5\% of responses (non-functioning) & 193 of 349 (55.3\%) \\
Items improving / worsening with repetition (FDR $q<0.05$) & 44 / 1 (of 164 testable) \\
\bottomrule
\end{tabular}

\vspace{2pt}
\parbox{\linewidth}{\footnotesize Difficulty and discrimination computed on first attempts; item--rest correlation uses each employee's accuracy on the remaining items (employees with $\geq$10 first attempts). Trend: per-item logistic regression of correctness on log$_2$(attempt) with employee-clustered SE, Benjamini--Hochberg correction. Distractors never chosen are not observable in the log and are not counted.}
\end{table}
